\subsection{The named completed action on a directly supplied tensor cell}
\label{sec:mpu_supplied_completed_local_action}

Let $W$ have physical size $d^2$ and supplied standard-form gates $u,v$.
Let $e:\{1,\ldots,b\}\to\{1,\ldots,d^2\}$ be a bijective relabelling,
and let $W'=e^*W$ be the actual reindexed tensor. The datum is explicitly
supplied, as in \cite{gap:mpu_qca_complete_block_heisenberg_scope}; no
preimage of $W$ under tensor blocking is introduced.

\begin{lemma}[Physical relabelling and complete window coordinates]
    \label{lem:mpu_supplied_relabel_coordinates}
    \lean{MPOTensor.IsMPU.reindexPhysical,
        MPOTensor.windowInputHeadEquiv,MPOTensor.windowOutputHeadEquiv,
        MPOTensor.twoSiteSuppliedCellEquiv,
        MPOTensor.twoSiteSuppliedInputWindowEquiv,
        MPOTensor.twoSiteSuppliedOutputWindowEquiv,
        MPOTensor.twoSiteSuppliedInputWindow,
        MPOTensor.twoSiteSuppliedOutputWindow}
    \notready
    \uses{def:mpu_supplied_window_coordinates}
    Physical relabelling preserves the actual periodic unitaries, since
    it simultaneously reindexes the row and column configurations.
    The singleton coordinate associated with a half-pair $c$ is the
    constant one-cell configuration $e^{-1}(\operatorname{pair}(c))$.
    The input head coordinates send $((j_0,j_2),c)$ to
    \begin{align}
        \bigl(e^{-1}(j_0),e^{-1}(\operatorname{pair}(c)),e^{-1}(j_2)\bigr).
    \end{align}
    The output coordinates send exterior halves $(a_L,a_R)$ and four
    output halves $(x_1,x_2,x_3,x_4)$ to
    \begin{align}
        \bigl(e^{-1}(\operatorname{pair}(a_L,x_1)),
        e^{-1}(\operatorname{pair}(x_2,x_3)),
        e^{-1}(\operatorname{pair}(x_4,a_R))\bigr).
    \end{align}
    Let $I_e(A)$ and $O_e(A)$ be the three-cell matrices obtained by
    adjoining identity exteriors to $A$ and $\Phi_{u,v}(A)$ in these
    complete coordinates.
\end{lemma}
\begin{proof}
    \notready
    Reindex each actual periodic matrix through the sitewise bijection;
    simultaneous reindexing preserves unitarity. The complete head
    equivalences already used by the finite-window calculation give the
    displayed coordinates, followed by sitewise physical relabelling.
\end{proof}

\begin{lemma}[Canonical buffered conjugation after physical relabelling]
    \label{lem:mpu_supplied_relabel_buffered_conjugation}
    \lean{MPOTensor.mpo_suppliedBufferedWindow_conjugation,
        MPOTensor.bufferedWindowMatrix_reindexPhysical,
        MPOTensor.mpo_reindexPhysical_suppliedBufferedWindow_conjugation,
        MPOTensor.twoSiteSuppliedInputWindow_eq_bufferedCell}
    \notready
    \uses{lem:mpu_supplied_relabel_coordinates,lem:mpu_direct_supplied_window_conjugation}
    Write $B_{p,s}(X)=I_p\otimes X\otimes I_s$ for actual canonical
    finite-window inclusion. Sitewise physical relabelling commutes with
    $B_{p,s}$. For every $p,s\ge0$, with $N=p+3+s$,
    \begin{align}
        \operatorname{MPO}_{W'}(N)\,B_{p,s}(I_e(A))\,
        \operatorname{MPO}_{W'}(N)^* & =B_{p,s}(O_e(A)).
    \end{align}
    Moreover $I_e(A)$ is exactly the singleton input matrix with one
    identity cell adjoined on each side.
\end{lemma}
\begin{proof}
    \notready
    The complete supplied-window coordinate formula equals canonical
    buffering by evaluating both matrices in the prefix/head/suffix
    configuration bijection. Apply the sitewise physical relabelling
    to the actual periodic conjugation identity. It commutes with
    concatenation of finite configurations and hence with identity
    buffering. The input head dictionary identifies its middle cell
    with the singleton coordinate, while the two exterior cells are
    identity factors.
\end{proof}

\begin{lemma}[Actual singleton output of the named automorphism]
    \label{lem:mpu_supplied_named_singleton_action}
    \lean{SpinChain.mpuQuasiLocalEquiv_supplied_singleton}
    \notready
    \uses{lem:mpu_supplied_relabel_buffered_conjugation,
        thm:mpu_actual_qca}
    Assume $b>0$. Let $\alpha_{W'}$ be the named actual bare-MPU
    automorphism and let $C_e(A)$ be the singleton-cell input matrix in
    the coordinates above. Then
    \begin{align}
        \alpha_{W'}\bigl(\iota_{[0,1)}(C_e(A))\bigr)
         & =\iota_{[-1,2)}(O_e(A)).
    \end{align}
    This is the actual supplied four-half-site output on cells $-1,0,1$.
    No compatible action, extra blocking scale or inverse-action datum
    is assumed.
\end{lemma}
\begin{proof}
    \notready
    The supplied standard form gives the MPU property of $W$, and
    physical reindexing gives that of $W'$. Choose exterior lengths
    $p=s=4D^4+1$. The resulting chain contains the singleton's full
    bare-MPU propagation buffer and the three-cell output window. The
    actual finite evaluation of the named automorphism therefore
    applies. Canonical local inclusions identify the singleton input
    with the central three-cell input window; the preceding periodic
    conjugation identity gives the stated output after faithful
    finite-region inclusion into the quasi-local algebra.
\end{proof}
